% ===================================================================
%  तक्ता क्रमांक १ — शाळा. — माध्यमिक शाळा. — वर्ग.
%
%  Ceci n'est PAS une transcription : c'est une traduction, faite en
%  2026, en marathi d'aujourd'hui. D'ou trois differences avec
%  texto/io et texto/fr, et trois seulement :
%
%   * pas de \nl ni de \cc — la traduction ne suit aucune ligne
%     imprimee, et les lignes de ce fichier ne sont que du confort ;
%   * pas de \begin{VUpage} — elle n'a ni feuillet ni folio ;
%   * le gras se pose sur le TERME seul, ponctuation dehors.
%
%  Les cles « %%K » sont celles de texto/io, macro pour macro pour
%  l'apparat, et les renvois \textsuperscript{(n)} visent les MEMES
%  objets de la planche, dans le MEME ordre.
%
%  POURQUOI CETTE COLONNE EXISTE, ALORS QUE « hi » EST DEJA LA. Parce
%  que le marathi n'est pas une graphie du hindi : c'est une autre
%  langue, qui s'ecrit dans le meme alphabet. Ce n'est meme pas le cas
%  du gourmoukhi et du shahmoukhi, qui sont une seule langue en deux
%  ecritures et qu'on traduit pourtant separement. Ici la grammaire
%  entiere differe :
%
%    आहे / आहेत       la copule    हिन्दी : है / हैं
%    चा / ची / चे     le genitif   हिन्दी : का / के / की
%    ला              le datif     हिन्दी : को
%    मध्ये / -त       le locatif   हिन्दी : में
%    नाही            la negation  हिन्दी : नहीं
%    हा / ही / हे     ce, cette    हिन्दी : यह
%    तो / ती / ते     celui-la     हिन्दी : वह
%    आणि             et           हिन्दी : और
%    पण              mais         हिन्दी : लेकिन
%    कारण            parce que    हिन्दी : क्योंकि
%    खूप             tres         हिन्दी : बहुत
%    काय / कसे / कुठे  quoi, comment, ou   हिन्दी : क्या / कैसे / कहाँ
%
%  ET L'APPARAT LUI-MEME CHANGE DE MOT, ce qui est rare : le marathi
%  dit तक्ता la ou le hindi dit तालिका, मालिका la ou il dit शृंखला,
%  et surtout प्रवेश pour une scene — le mot de son theatre — la ou
%  le hindi dit दृश्य. Les trois motifs de html.py ont donc chacun
%  gagne un membre.
%
%  LA CLASSE MARATHIE A SON VOCABULAIRE ENTIER, ET IL N'EST PAS CELUI
%  DU HINDI :
%
%    फळा (1)        le tableau noir
%    खडू (3)        la craie
%    दप्तर (33, 57) le cartable
%    वही (28)       le cahier
%    खरडवही (35)    le cahier de brouillon
%    खोडरबर (36)    la gomme
%    पाटी (42)      l'ardoise
%    शाईदाणी (25)   l'encrier
%    टीपकागद (41)   le buvard
%    कपाट (15)      l'armoire
%    बाक (32)       le banc
%    झरोका (78)     le vasistas
%    मधली सुट्टी (76) la recreation
%    पटांगण (75)    la cour de l'ecole
%    गुरुजी          le maitre
%    मुख्याध्यापक (71) le proviseur
%    शिपाई (77)      le concierge
%    हुशार (19)      bon eleve
%
%  « शिपाई » EST LE MEME EMPLOI QUE LE فرّاش EGYPTIEN : l'homme de
%  service de l'ecole, celui qui porte le registre des absents. Deux
%  colonnes ecrites de suite, deux ecoles reelles, et la meme
%  fonction nommee sans periphrase.
%
%  « postigo » (t01, renvoi 78) EST UN VASISTAS, ET LA PLANCHE LE DIT.
%  Verifie sur la gravure en 2026 : dans la rangee HAUTE de la
%  verriere, deux petits vantaux sont graves battants, gonds sur un
%  montant vertical, ouverts vers l'interieur. C'est le seul ouvrant,
%  et il ne sert qu'a aerer — donc झरोका, ni lucarne ni volet ni
%  meneau.
% ===================================================================

%%K t01-apar-0 apar
\VUsaut{2.0mm}
\VUcentre{12.6pt}{120}{स्पष्टीकरण-पुस्तिका}
\VUsaut{3.2mm}
\VUcentre{8.4pt}{160}{सोबतची}
\VUsaut{2.6mm}
\VUtitre{15.6pt}{0}{देल्मास साहाय्यक तक्त्यांची}
\VUsaut{3.4mm}
\VUornamento{9pt}
\VUsaut{5.0mm}
\VUcentre{13.2pt}{90}{पहिली मालिका}
\VUsaut{3.0mm}
\VUfilet{16mm}
\VUsaut{5.6mm}

%%K t01-apar-1 apar
\VUtitre{15.0pt}{60}{तक्ता क्रमांक १}
\VUsaut{1.4mm}
\VUtitre{11.4pt}{0}{शाळा. --- माध्यमिक शाळा. --- वर्ग.}
\VUsaut{2.2mm}
\VUfilet{20mm}
\VUsaut{6.4mm}

%%K t01-tit-1 sub
\VUpk{10.2pt}{40}{सर्वसाधारण वर्णन.}
\VUsaut{3.0mm}

%%K t01-01-1 p
१. --- हा तक्ता माध्यमिक \VUgras{शाळेतला} एक \VUgras{वर्ग} दाखवतो. या वर्गात आठ
\VUgras{टेबले} \textsuperscript{(18)} आणि आठ \VUgras{बाक} \textsuperscript{(32)}
आहेत. प्रत्येक बाकावर तीन विद्यार्थी बसतात. आणि \VUgras{गुरुजी}
\textsuperscript{(7)} आपल्या \VUgras{खुर्चीवर} \textsuperscript{(8)} आपल्या
\VUgras{व्यासपीठासमोर} \textsuperscript{(6)} बसले आहेत.

%%K t01-02-1 p
२. --- व्यासपीठाच्या डावीकडे काचेचे \VUgras{कपाट} \textsuperscript{(15)} आहे.
उजवीकडे \VUgras{फळा} \textsuperscript{(1)} आहे, त्याच्यासोबत \VUgras{फडके}
\textsuperscript{(2)} आणि \VUgras{खडू} \textsuperscript{(3)}.

%%K t01-02-2 p
आणि व्यासपीठाच्या वर \VUgras{भिंतीवरचे तक्ते} \textsuperscript{(9, 11, 12)} व
\VUgras{नकाशा} \textsuperscript{(10)} टांगलेले आहेत.

%%K t01-03-1 p
३. --- सगळे विद्यार्थी आपापल्या \VUgras{जागेवर} \textsuperscript{(17)} नाहीत. जॉन
\textsuperscript{(4)} फळ्याजवळ उभा आहे; तो फडके धरून \VUgras{भूमितीच्या आकृत्या}
पुसतो आहे.

%%K t01-03-2 p
आणि पीटर व्यासपीठाजवळ आहे; तो \VUgras{लेखी कामे} \textsuperscript{(5)} गोळा करून
गुरुजींकडे नेतो आहे.

%%K t01-04-1 p
४. --- \VUgras{हे} गुरुजी \VUgras{आधुनिक भाषांचे} शिक्षक आहेत. ते विद्यार्थ्यांना
परकीय भाषा बोलायला, वाचायला आणि लिहायला शिकवतात \VUgras{(जर्मन, इंग्रजी, स्पॅनिश,
इटालियन, रशियन} किंवा \VUgras{फ्रेंच)}.

%%K t01-05-1 p
५. --- वर्ग खूप प्रशस्त आणि खूप उजेडाचा आहे. त्याच्या टोकाला रुंद \VUgras{खिडकी}
आहे, तिला दोन \VUgras{झरोके} \textsuperscript{(78)} आहेत, आणि \VUgras{काचेचे दार}
आहे --- हवा आणि उजेड येण्यासाठी.

%%K t01-05-2 p
आणि हा वर्ग थंड नाही. त्याच्या मधोमध, ऊब मिळावी म्हणून, एक खूप चांगली \VUgras{शेगडी}
\textsuperscript{(46)} तिच्या \VUgras{धुराड्यासह} \textsuperscript{(45)} आहे.

%%K t01-05-3 p
आणि आडभिंतीवर \VUgras{खुंट्या} \textsuperscript{(58)} बसवलेल्या आहेत; त्यांवर
विद्यार्थी आपल्या टोप्या व कोट टांगतात.

%%K t01-tit-2 sub
\VUsaut{5.2mm}
\VUpk{10.2pt}{40}{खेळाचे पटांगण.}
\VUsaut{3.6mm}

%%K t01-06-1 p
६. --- \VUgras{घड्याळ} \textsuperscript{(63)} नऊ वाजून पाच मिनिटे दाखवते आहे.

%%K t01-06-2 p
\VUgras{मधली सुट्टी} \textsuperscript{(76)} आत्ताच संपली आणि तास पुन्हा सुरू होतो
आहे. खेळ \VUgras{पटांगणात} \textsuperscript{(75)} होतो. आपल्याला काही विद्यार्थी
खेळताना दिसतात. \VUgras{शिक्षकांचा मदतनीस} \textsuperscript{(74)} त्यांच्यावर लक्ष
ठेवतो आहे.

%%K t01-07-1 p
७. --- पटांगणात आणखीही निरनिराळी माणसे दिसतात, म्हणजे \VUgras{निरीक्षक}
\textsuperscript{(73)}, \VUgras{मुख्याध्यापक} \textsuperscript{(71)}, आणि
\VUgras{उपमुख्याध्यापक} \textsuperscript{(72)}.

%%K t01-07-2 p
निरीक्षक शाळांची तपासणी करतो, मुख्याध्यापक शाळा चालवतात, आणि उपमुख्याध्यापक
अभ्यासावर लक्ष ठेवतात.

%%K t01-07-3 p
आणि चौथी व्यक्ती म्हणजे \VUgras{शिपाई} \textsuperscript{(77)}. गैरहजर
विद्यार्थ्यांची नोंद करण्यासाठी तो काखेत हजेरीपुस्तक घेऊन फिरतो.

%%K t01-08-1 p
८. --- पटांगणाच्या टोकाला \VUgras{व्यायामाची साधने} \textsuperscript{(70)} आहेत,
आणि \VUgras{हस्तकलेची} \textsuperscript{(69)} कार्यशाळा आहे.

%%K t01-08-2 p
आणि तक्त्याच्या उजवीकडे \VUgras{संगीताचा} व \VUgras{चित्रकलेचा} \VUgras{वर्ग}
दिसू लागतो.

%%K t01-noto-1 noto
\VUnotes{143.2mm}{%
(*) \textit{बाप्तिस्म्याच्या नावांसाठीही} ती लॅटिन रूपात लिहावीत असा सल्ला
दिला जातो. असंख्य फरक टाळण्याचा तोच एक मार्ग आहे. शिवाय \textit{Giovanni,
Jean, Johann, Jan, Hans, John, Ivan,} इत्यादींपैकी निवड कशी करायची?
बाप्तिस्म्याच्या नावांसाठी वेगळा शब्दकोश करायचा का? लॅटिनमधून त्यांचे प्रथमा
रूप घेऊन म्हणू या: \VUgras{Ioannes, Ioanna; Iulius, Iulia; Iakobus;}
\VUgras{Andreas; Lukas.} (Gram. Kompleta).%
}

%%K t01-tit-3 sub
\VUpk{10.2pt}{40}{सविस्तर वर्णन.}
\VUsaut{2.4mm}

%%K t01-tit-4 sub
\VUpk{10.2pt}{40}{हुशार विद्यार्थी.}
\VUsaut{3.0mm}

%%K t01-09-1 p
९. --- व्यासपीठाच्या उजवीकडच्या पहिल्या बाकावरचे विद्यार्थी म्हणजे
\VUgras{हेन्री} (*), \VUgras{स्टीफन} आणि \VUgras{चार्ल्स}.

%%K t01-09-2 p
हेन्री खूप लक्ष देतो आणि कष्टाळू आहे; तो गुरुजींचे ऐकतो.

त्याच्या टेबलावर \VUgras{वही} \textsuperscript{(28)}, \VUgras{पुस्तक}
\textsuperscript{(29)}, \VUgras{शब्दकोश} \textsuperscript{(30)} आणि
\VUgras{कंपासपेटी} \textsuperscript{(31)} आहे. त्याच्या पुस्तकाला पुष्कळ
\VUgras{पाने} आहेत, प्रत्येक प्रकरणात अनेक \VUgras{परिच्छेद} आहेत आणि प्रत्येक
\VUgras{ओळीत} पुष्कळ \VUgras{शब्द} आहेत.

%%K t01-09-4 p
त्याचा शेजारी स्टीफन \textsuperscript{(24)} मेहनती आहे. पुढच्या वेळचा धडा तो
आपल्या वहीत लिहून घेतो आहे. त्याचे \VUgras{हस्ताक्षर} सुंदर आहे. त्याचे पुस्तक
मिटलेले आहे. आपल्याला फक्त \VUgras{मुखपृष्ठ} \textsuperscript{(27)} दिसते.
त्याच्याजवळ त्याचा \VUgras{कंपास} \textsuperscript{(26)} आहे.

%%K t01-09-5 p
चार्ल्स \textsuperscript{(23)} हाताने आपले पुस्तक धरून आहे; धडा पुन्हा वाचण्यासाठी
तो ते उघडतो आहे. प्रत्येक विद्यार्थ्याप्रमाणे त्याच्याही समोर \VUgras{शाईदाणी}
\textsuperscript{(25)} आहे. तो आज्ञाधारक आहे.

%%K t01-10-1 p
१०. --- शेजारच्या टेबलावर \VUgras{लुई}, \VUgras{अँटनी} आणि \VUgras{पॉल} आहेत. लुई
आपला \VUgras{धडा} \textsuperscript{(22)} म्हणून दाखवतो आहे आणि गुरुजींच्या
\VUgras{प्रश्नांची} उत्तरे देतो आहे. अँटनी \textsuperscript{(21)} अजिबात लक्ष देत
नाही, कधीकधी तर गुरुजी त्याला शिक्षा करतात. पॉल \textsuperscript{(20)} चांगला
सोबती आहे, प्रेमळ आणि मदत करणारा. आणि तो खूप लक्ष देतो.

%%K t01-tit-5 sub
\VUsaut{4.2mm}
\VUpk{10.2pt}{40}{वाईट विद्यार्थी.}
\VUsaut{3.2mm}

%%K t01-11-1 p
११. --- हेन्रीच्या मागे \VUgras{जॉर्ज} \textsuperscript{(38)} बसतो. तो वाईट मुलगा
नाही, पण \VUgras{बडबड्या}, बेलक्ष आणि खोडकर आहे. त्याच्या \VUgras{चंचलपणामुळे} व
\VUgras{अवखळपणामुळे} तासात \VUgras{गोंधळ} होतो आणि कितीदा तरी त्याला कडक
\VUgras{ताकीद} मिळते. त्याची \VUgras{खरडवही} \textsuperscript{(35)} मळकट आहे, तो
आपल्या \VUgras{टाकाने} \textsuperscript{(34)} तिच्यावर \VUgras{शाईचे डाग पाडतो}, आणि
म्हणूनच तो नेहमी आपले \VUgras{खोडरबर} \textsuperscript{(36)} वापरतो; त्याचे
\VUgras{दप्तर} \textsuperscript{(33)} फाटलेले आहे, कारण तो खूप निष्काळजी आहे. मग
\VUgras{आधुनिक भाषांच्या} वर्गात तो आपली \VUgras{पेन्सिलपेटी} \textsuperscript{(37)}
कशाला आणतो?

%%K t01-11-2 p
त्याचा शेजारी एडमंड \textsuperscript{(39)} त्याला आवडत नाही. खरे तर तो जरा आळशी
आहे, पण गप्प आणि शांत आहे. तो बडबड करत नाही; मात्र ऐकण्याऐवजी त्याला आपल्या
\VUgras{वाचनाच्या पुस्तकातल्या} \VUgras{गोष्टी} वाचायला आवडतात.

%%K t01-11-3 p
या बाकावरचा तिसरा विद्यार्थी म्हणजे \VUgras{लुई} \textsuperscript{(40)}. तो मेहनती
आहे. तो पांढऱ्या \VUgras{कागदावर} लिहितो आहे. आणि समोर त्याने आपला
\VUgras{टीपकागद} \textsuperscript{(41)} ठेवला आहे.

%%K t01-12-1 p
१२. --- गुरुजींच्या डावीकडच्या दुसऱ्या बाकावरचे विद्यार्थी खूप लहान आहेत.
पहिला, लहानगा \VUgras{एमिल}, अजून नीट लिहू शकत नाही; तो आपली \VUgras{पाटी}
\textsuperscript{(42)}, आपली \VUgras{पाटीपेन्सिल} आणि आपला \VUgras{स्पंज}
\textsuperscript{(43)} आणतो.

%%K t01-12-2 p
त्याच्या शेजारी \VUgras{ऑगस्ट} आणि \VUgras{बर्नार्ड} \textsuperscript{(44)} आहेत.
हा चांगला अभ्यास करतो, पण त्याचे \VUgras{कपडे} फार अव्यवस्थित असतात.

%%K t01-13-1 p
१३. --- त्यांच्या मागे तीन \VUgras{आळशी} आहेत. \VUgras{अल्बर्ट} आपले धडे शिकत
नाही. \VUgras{जेम्स} \textsuperscript{(47)} ऐकण्याऐवजी झोपतो. त्याच्या
\VUgras{आळसामुळे} त्याला \VUgras{बोलणी} आणि \VUgras{शिक्षाही} मिळते. आणि शेवटी
\VUgras{ख्रिश्चन} \textsuperscript{(48)} फारच बेजबाबदार आहे. त्याची
\VUgras{वागणूक} वाईट आहे आणि सुधारण्याचा तो अजिबात प्रयत्न करत नाही.

%%K t01-14-1 p
१४. --- याउलट त्याचे शेजारी \VUgras{सुस्वभावी} आहेत. \VUgras{अर्नेस्ट}
\textsuperscript{(51)} याला टापटीप आणि नियमितपणा आवडतो; तो नेहमी आपली
\VUgras{पट्टी} \textsuperscript{(50)} आणि आपली \VUgras{पेन्सिल} \textsuperscript{(49)}
वापरतो. तो \VUgras{बाहेरून येणारा विद्यार्थी} आहे.

%%K t01-14-2 p
\VUgras{फ्रान्सिस} \textsuperscript{(52)} \VUgras{वसतिगृहात राहणारा} आहे; तो गणवेश
घालतो, शाळेतच जेवतो आणि झोपतो. तो चांगला \VUgras{सोबती} आहे.

%%K t01-14-3 p
मॉरिस आपल्या \VUgras{चाकूने} \textsuperscript{(56)} आपली \VUgras{पेन्सिल} तासतो; तो
खूप काळजी घेणारा आहे.

%%K t01-14-4 p
तो आपली सगळी पुस्तके आणतो, आपले \VUgras{व्याकरण} \textsuperscript{(55)}, आपली
\VUgras{टिपणवही} \textsuperscript{(54)}, आणि \VUgras{लिहिण्याची उशीसुद्धा}
\textsuperscript{(53)}. त्याचे \VUgras{दप्तर} \textsuperscript{(57)} त्याच्या मागे
जमिनीवर आहे.

%%K t01-14-5 p
वर्गाच्या टोकाची दोन टेबले \VUgras{हुशार विद्यार्थ्यांनी} \textsuperscript{(19)}
भरलेली आहेत.

%%K t01-tit-6 sub
\VUsaut{4.6mm}
\VUpk{10.2pt}{40}{चित्रकला आणि संगीत.}
\VUsaut{3.4mm}

%%K t01-15-1 p
१५. --- \VUgras{चित्रकलेच्या वर्गात} भिंतीवर सुंदर \VUgras{नमुने} टांगलेले आहेत
आणि छताला \VUgras{शेडसह} एक \VUgras{दिवा} \textsuperscript{(62)} लटकवलेला आहे.
\VUgras{शिक्षक} \textsuperscript{(60)} खूप संयमी आहेत; ते विद्यार्थ्यांची
\VUgras{चित्रे} \textsuperscript{(61)} सुधारतात आणि त्यांना निसर्गावरून
\VUgras{अर्धपुतळा} \textsuperscript{(59)} काढायला शिकवतात.

%%K t01-16-1 p
१६. --- \VUgras{संगीताचा वर्ग} फार रंजक वाटतो. तिथे \VUgras{संगीत} वाचायला
शिकतात, \VUgras{रेघांच्या ओळीवर} लिहिलेले \VUgras{स्वर} \textsuperscript{(67)}
म्हणायला शिकतात (do, re, mi, fa, sol, la, si\,=\,C. D. E. F. G. A. B.),
\VUgras{किल्ल्या} ओळखायला शिकतात आणि गोड \VUgras{चाली} गायला शिकतात. संगीताची पाने
\VUgras{मेजपाटीवर} \textsuperscript{(65)} ठेवतात. एक विद्यार्थी \VUgras{पियानो
वाजवतो} \textsuperscript{(64)} आणि \VUgras{संगीतशिक्षक} \textsuperscript{(66)}
\VUgras{ताल धरतात} \textsuperscript{(68)}.

%%K t01-17-1 p
१७. --- आपल्याला \VUgras{हस्तकलेच्या} \textsuperscript{(69)}
\VUgras{कार्यशाळेचा} एक कोपरा दिसतो, जिथे मुले लाकडाला रंधा मारायला आणि लोखंडाला
कानस लावायला शिकू शकतात. हे सगळ्या माध्यमिक शाळांत नसते.

%%K t01-tit-7 sub
\VUsaut{5.0mm}
\VUpk{10.2pt}{40}{विज्ञानाचा वर्ग.}
\VUsaut{3.6mm}

%%K t01-18-1 p
१८. --- गणिताचे शिक्षक विद्यार्थ्यांना \VUgras{भूमिती, अंकगणित, गणिती क्रिया}
आणि \VUgras{मापनपद्धती} शिकवतात. तक्त्यावर आपल्याला मुख्य भूमितीय आकृत्या
दिसतात: \VUgras{वर्तुळ} \textsuperscript{(a)}, \VUgras{चौरस} \textsuperscript{(b)},
\VUgras{त्रिकोण} \textsuperscript{(c)}, \VUgras{रेषा} \textsuperscript{(d)}.

%%K t01-18-2 p
आपल्याला एक \VUgras{क्रियाही} दिसते; ती \VUgras{बेरीज} नाही, \VUgras{वजाबाकी}
नाही, \VUgras{गुणाकारही} नाही --- ती \VUgras{भागाकार} \textsuperscript{(e)} आहे.

%%K t01-19-1 p
१९. --- मापनपद्धतीच्या तक्त्यावर आपल्याला दिसते: \VUgras{मीटर}
\textsuperscript{(a)}, \VUgras{तराजू} \textsuperscript{(b)} आणि त्याची
\VUgras{वजने}, \VUgras{लिटर} \textsuperscript{(e)}, \VUgras{डेकालिटर}
\textsuperscript{(f)}, \VUgras{डेसिलिटर} आणि \VUgras{घनमीटर} \textsuperscript{(d)}.

%%K t01-19-2 p
विद्यार्थी \VUgras{निसर्गविज्ञानही} \textsuperscript{(11)} शिकतात --- प्राणिशास्त्र
\textsuperscript{(a)}, वनस्पतिशास्त्र \textsuperscript{(b)}, भूशास्त्र
\textsuperscript{(c)} --- आणि \VUgras{इतिहास} \textsuperscript{(12)}.

%%K t01-19-3 p
कपाटात \VUgras{भौतिकशास्त्राची} \textsuperscript{(15)} आणि
\VUgras{रसायनशास्त्राची} \textsuperscript{(16)} उपकरणे आहेत.

%%K t01-19-4 p
कपाटावर आहेत: \VUgras{गोल} \textsuperscript{(14)}, \VUgras{शंकू} आणि
\VUgras{पृथ्वीगोल} \textsuperscript{(13)}.

%%K t01-19-5 p
तक्त्यांच्या वर जगाचे पाच भाग दाखवणारा \VUgras{नकाशा} टांगलेला आहे:
\VUgras{युरोप} \textsuperscript{(g)}, \VUgras{आशिया} \textsuperscript{(e)},
\VUgras{आफ्रिका} \textsuperscript{(f)}, \VUgras{अमेरिका} \textsuperscript{(ab)} आणि
\VUgras{ऑस्ट्रेलिया} \textsuperscript{(h)}, आणि दोन मोठे महासागर ---
\VUgras{पॅसिफिक} \textsuperscript{(c)} आणि \VUgras{अटलांटिक} \textsuperscript{(d)}.

\VUsaut{6.0mm}
\VUfilet{22mm}
